\subsection{Supuestos y Restricciones [Pendientes: 1. actualizar disponibilidad y preparación de datos, 2. comprobar recursos y condiciones de aplicación]}\label{sec:spmp-6-3}
La viabilidad del trabajo presupone que será posible acceder a datos territoriales utilizables, configurar una revisión de WFDS adecuada al problema y disponer de capacidad de cómputo y almacenamiento proporcional al conjunto de escenarios que se acuerde. También presupone la posibilidad de recibir acompañamiento académico y de conversar con personas que representen al usuario previsto. Cada condición se conservará como supuesto hasta contar con evidencia suficiente, de manera que una fuente identificada, una oferta de infraestructura o una intención de participación no se registren como disponibilidad confirmada. En particular, la información de vegetación existente requiere estudio de estructura, propiedades y condiciones de transferencia antes de convertirse en entradas del simulador~\cite{datos2026}.

Las restricciones vigentes son el producto de escritorio, el sector territorial acotado, el calendario académico y la capacidad efectivamente disponible para generar y conservar resultados. El clima consultado mediante una API introduce una dependencia que impide equiparar precálculo con funcionamiento completo sin conexión, mientras que el muestreo por amenaza y la agregación requieren comprobaciones propias. El inventario de disponibilidad deberá actualizarse conforme se consigan archivos, se preparen datos y se validen entradas, conservando fuente, fecha, unidad, cobertura y procedimiento de transformación. Esa actualización alimentará la gestión de riesgos y las estimaciones, cuyos tratamientos se desarrollan respectivamente en las secciones \ref{sec:spmp-12-2} y \ref{sec:spmp-9-1}, sin declarar aprobados parámetros científicos todavía en estudio.
